% ICLR preliminaries excerpt; requires amsmath and amssymb.
% Citation: section1_references.bib. Review: preliminaries_review.md.
% Replaces the former standalone section1_single_mode.tex motivation.
% Outline: behavior-cloning setup -> MSE baseline -> multimodality explanation
% -> evidence against that explanation -> successful regression -> M-estimation
% -> heteroscedastic regression -> demonstration-data analysis.

\section{Preliminaries}
\label{sec:preliminaries}

% Setup: define the task without restricting the policy class.
We study behavior cloning from demonstrations $\mathcal{D}=\{(o_i,a_i)\}$,
where $o_i$ contains the observations and task instruction, and $a_i$ is an
action chunk. The goal is to learn a policy that generates action chunks
conditioned on these inputs to accomplish the demonstrated tasks.
A standard baseline is direct action regression~\citep{robomimic2021}, which
trains a network $f_\theta(o)$ to predict an action chunk by minimizing the
mean squared error (MSE):
\begin{equation}
    \mathcal{L}_{\mathrm{MSE}}(\theta)
    = \mathbb{E}_{(o,a)\sim\mathcal{D}}
      \bigl[\lVert f_\theta(o)-a\rVert_2^2\bigr].
    \label{eq:prelim-mse}
\end{equation}

% Background: frame diffusion and flow together as denoising-based policy learning.
However, denoising-based policies have outperformed conventional
behavior-cloning baselines on challenging manipulation
tasks~\citep{chi2023diffusionpolicy,pan2026much}.
A common explanation attributes this advantage to multimodal demonstrations:
generative policies can represent multiple valid actions for the same
observation, whereas direct MSE regression can average incompatible
actions~\citep{chi2023diffusionpolicy}.

\subsection{Multimodality and deterministic policies}
\label{sec:prelim-regression-gap}

% Evidence against the explanation: unimodal samples and deterministic data.
\citet{pan2026much} challenge this explanation. In their Flow-policy probes,
action samples form a single cluster even at symmetry-critical or ambiguous
states, and executing the mean sampled action has little effect on task
success. The Flow--regression gap also persists on data collected from a
deterministic expert, indicating that multimodality does not explain
the performance gap.

% Evidence and takeaway: regression can succeed without multimodal sampling.
Deterministic regression can match Flow's performance. The minimal iterative
policy (MIP) of \citet{pan2026much}, trained with squared-error supervision
and noise injection, generates actions deterministically in two network
evaluations.
Across seven challenging tasks, the authors report an average success rate
relative to Flow of $1.02$ for MIP and $0.74$ for direct regression
(their Figure~1). These results show that modeling multiple action modes
is not necessary for strong control performance: regression-based policies
can match Flow.

\subsection{M-estimation and robust regression}
\label{sec:prelim-m-estimation}

% Background: connect residual penalties, maximum likelihood, and MSE.
M-estimation generalizes least squares by replacing the squared error with
a residual penalty~\citep{huber1964robust}. For the continuous action
channels, let $r_\theta=f_\theta(o)-a\in\mathbb{R}^d$ denote the flattened
action-chunk residual. An M-estimator minimizes
\begin{equation}
    \mathcal{L}_{\rho}(\theta)
    = \mathbb{E}_{(o,a)\sim\mathcal{D}}
      \bigl[\rho(r_\theta)\bigr].
    \label{eq:prelim-m-estimation}
\end{equation}
Choosing $\rho(r)=-\log q(r)$ for a residual density $q$ gives maximum
likelihood estimation (MLE). In particular, an isotropic Gaussian
$q=\mathcal{N}(0,\sigma^2 I_d)$ with fixed $\sigma>0$ yields the MSE
objective in Eq.~\ref{eq:prelim-mse}, up to a positive multiplicative
factor and an additive constant.

% Mechanism: the residual derivative explains sensitivity to large errors.
Robust regression reduces sensitivity to extreme residuals through the
shape of $\rho$. The prediction gradient $\psi(r)=\nabla_r\rho(r)$ equals
$2r$ for squared error and grows linearly with residual magnitude.
By contrast, coordinate-wise Huber penalties are quadratic near zero and
linear beyond a threshold, yielding bounded derivatives~\citep{huber1964robust}.

\subsection{Heteroscedastic regression}
\label{sec:prelim-heteroscedastic}

% Background: input-dependent scale leads to a normalized likelihood.
Heteroscedastic regression allows the conditional residual scale to vary
with the input. A common formulation jointly predicts an action mean
$f_\theta(o)$ and a positive scalar scale $\sigma_\theta(o)$, using the
Gaussian model
$p_\theta(a\mid o)=\mathcal{N}\!\bigl(a;f_\theta(o),
\sigma_\theta(o)^2 I_d\bigr)$~\citep{nix1994estimating,seitzer2022pitfalls}.
Omitting additive constants, its negative log-likelihood is
\begin{equation}
    \mathcal{L}_{\mathrm{HG}}(\theta)
    = \mathbb{E}_{(o,a)\sim\mathcal{D}}
      \left[
        \frac{\lVert r_\theta\rVert_2^2}{2\sigma_\theta(o)^2}
        + d\log\sigma_\theta(o)
      \right].
    \label{eq:prelim-hg}
\end{equation}
Inverse-variance weighting measures each residual relative to its predicted
local scale, while the logarithmic term penalizes scale inflation.
This adapts the scale across inputs but retains a quadratic penalty on
the standardized residual $r_\theta/\sigma_\theta(o)$.

% Synthesis and transition: distinguish scale from tails before probing data.
Together, these formulations distinguish input-dependent scale from
sensitivity to extreme residuals, while retaining deterministic inference
through $f_\theta(o)$. We next examine both properties in the demonstration
data.
